


This continuation proves the continuous-phase version of H29--H31
and the exact maps to the completed targets in H64 of the delivered
\emph{All-holonomy descent of the original theta quotient}.
The delivered NOTE.tex was read completely, all 1152 lines;
its SHA-256 is
\begin{center}\small
\texttt{444f82a8fa420e6f471ef448442085f3a93acb137069e1505d155faab051b881}.
\end{center}
The archive SHA-256 is
\begin{center}\small
\texttt{f0506a27cd4c9012f43fa70c3d1f8f996bbb6db01b0eea2944b34dcf290ca8f0}.
\end{center}
All source coordinates, masses, phases, relation columns, and full
primary multiplicities below are retained. No finite test or Lean
execution is used as a proof.
The sampled-density argument uses the full analytic proof
\emph{Periodized source density review}, PSD8--PSD14,
whose TeX SHA-256 is
\begin{center}\small
\texttt{51102189bcc7a45610464143a5c0bee21c08059c317024e6e856bc7de993fe74}.
\end{center}
The completed-fibre calculation develops H64 and is also proved
in the companion completed-quotient and holonomy-intake derivations.
Its repetition here supplies the maps needed for this continuous
common-image complex, including the mean-zero contraction and
the precise weighted relation gap.

\section{Original coefficient spaces and actual phase norms}

Fix the original monic polynomial $\chi$, its full degree $q\ge1$,
and an integer $N\ge q-1$. Put
\[
\begin{gathered}
 E=\C[S]/(\chi),\quad P_N=\C[S]_{\le N},\quad
 J=J_N:P_N\longrightarrow E,\\
 B_N=\ker J=\chi P_{N-q},\quad A=M_S:E\longrightarrow E,
\end{gathered}
 \tag{HCD1}
\]
where $P_{-1}=0$. Use the original increasing-power coefficient
bases. Let $B:\C^r\to P_N$, $r=N-q+1$, be multiplication by
the literal $\chi$, so that $\im B=B_N$. At $r=0$ it is the
unique map from the zero coefficient space.

Let $\Theta=[0,2\pi)$ carry its actual Fourier phase measure
$d\mu(\theta)=d\theta/(2\pi)$. Retain the measurable Hermitian
phase source matrices $M_\theta$ and their mean
\[
 M=\int_\Theta M_\theta\,d\mu,\qquad
 0<a I_{P_N}\preceq M_\theta\preceq b I_{P_N}<\infty
 \quad\text{almost everywhere},\qquad a,b>0.
 \tag{HCD2}
\]
The original source isometry H13--H17 identifies $M$ with the
original source Gram, including its full constant mass.
For the actual source, the bounds in (HCD2) are obtained from
H39--H41, also proved in HG3, as follows. Let
$\Psi_N:P_N\to L^2(\R;\mathcal K)$ be the original finite
source columns after the half-density map, with
$\mathcal K=L^2(\R^{k-1},dz)$. Fix $\beta>0$ for the retained
weighted source Grams, whose finite entries are supplied by
the original source regularity:
\[
 M_{\beta,\pm}=(e^{\pm\beta r}\Psi_N)^*
                    (e^{\pm\beta r}\Psi_N),\qquad
 \kappa_{\beta,N}=
 \max_{v\ne0}\frac{v^*(M_{\beta,+}+M_{\beta,-})v}{v^*Mv}.
 \tag{HCD2a}
\]
These are the original weights $x_k^{\pm2\beta}$; all original
source measures remain present. The maximum is finite and positive
because $M>0$ and the two weighted Grams are finite positive forms.
For $\mathcal C(s)=\int\Psi_N(r)^*\Psi_N(r+s)\,dr$ and $s\ge0$,
write the integrand at $v$ as the inner product of
$e^{-\beta r}\Psi_N(r)v$ and
$e^{\beta(r+s)}\Psi_N(r+s)v$, multiplied by $e^{-\beta s}$.
Cauchy--Schwarz, followed by translation in the second integral,
gives
\[
 |v^*\mathcal C(s)v|
 \le e^{-\beta s}
 \sqrt{(v^*M_{\beta,-}v)(v^*M_{\beta,+}v)} .
\]
For $-s$ the same bound follows from
$\mathcal C(-s)=\mathcal C(s)^*$.
Unfolding the literal Zak sum gives
$M_\theta=\sum_{j\in\mathbb Z}e^{ij\theta}\mathcal C(jL)$.
The correlation bound and polarization give absolute convergence
of each matrix entry, uniformly in phase. Pairing $j$ and $-j$,
using $2\sqrt{uv}\le u+v$, and summing the geometric series prove
\[
 -\delta_LM\preceq M_\theta-M\preceq\delta_LM,\qquad
 \delta_L=\frac{\kappa_{\beta,N}}{e^{\beta L}-1}.
 \tag{HCD2b}
\]
For any chosen $0<\eta<1$, the explicit finite circumference
\[
 L\ge\frac1\beta\log(1+\kappa_{\beta,N}/\eta)
 \quad\Longrightarrow\quad
 \delta_L\le\eta<1
 \tag{HCD2c}
\]
therefore supplies (HCD2), with
$a=(1-\delta_L)\lambda_{\min}(M)$ and
$b=(1+\delta_L)\lambda_{\max}(M)$ in the original coefficient
frame. The phase mean is still exactly $M$: every nonzero
Fourier mode integrates to zero against $d\theta/(2\pi)$.
The concrete choice $\eta=1/2$ gives
$L=\beta^{-1}\log(1+2\kappa_{\beta,N})$ and $\delta_L=1/2$.
These are actual finite-source bounds at the stated cutoff;
no bound uniform in $N$ or the arithmetic packet is asserted.
The argument below also applies to any fixed $L$ for which
(HCD2) is already established.

Define the Hilbert space
\[
 \mathscr P=L^2(\Theta;P_N,M_\theta\,d\mu),\qquad
 \|P\|_{\mathscr P}^2=\int_\Theta P(\theta)^*M_\theta
 P(\theta)\,d\mu.
\]
By (HCD2), its elements are exactly coefficientwise $L^2$
functions, with equivalent norms. Equality means equality
almost everywhere. The coefficientwise Bochner integral
$\av P=\int P\,d\mu$ exists, since coefficientwise $L^2$
implies $L^1$ on this finite phase measure. Every fixed
coefficient map commutes with this integral, entry by entry.

The closed relation space and common-image space are
\[
 \mathscr B=L^2(\Theta;B_N),\qquad
 \mathscr P^{\rm eq}
 =\{P\in\mathscr P:JP(\theta)=x\text{ a.e. for a fixed }x\in E\}.
 \tag{HCD3}
\]
They carry the restricted source norms. They are closed:
$J$ is continuous on coefficientwise $L^2$, $B_N=\ker J$
is closed, and constant $E$-valued functions form a closed
subspace. For the last assertion, averaging followed by
repetition is a bounded idempotent with precisely that image.

\section{The continuous cochain splitting and its full contraction}

Use complexes in degrees $0,1$ with actual inclusion differentials:
\[
 \mathcal C_N=[B_N\hookrightarrow P_N],\qquad
 \mathscr C_N=[\mathscr B\hookrightarrow\mathscr P^{\rm eq}].
 \tag{HCD4}
\]
Repetition $\Delta v(\theta)=v$ and averaging $\av$ are cochain
maps. In degree one, $J\av P=\int JP\,d\mu=x$; in degree
zero the mean belongs to $B_N$. The equality $\av\Delta=I$
uses the literal phase measure. Repetition is an isometry:
\[
 \|\Delta v\|_{\mathscr P}^2
 =v^*\left(\int M_\theta\,d\mu\right)v=v^*Mv.
\]
Averaging is bounded with the explicit estimate
\[
 \|\av P\|_M^2
 \le b\|\av P\|_{\rm coeff}^2
 \le b\int\|P(\theta)\|_{\rm coeff}^2\,d\mu
 \le (b/a)\|P\|_{\mathscr P}^2.
 \tag{HCD5}
\]
This uses scalar Cauchy--Schwarz for the coefficient vector
and does not replace the arithmetic mass by one.

Put $\mathscr B^0=\{b\in\mathscr B:\av b=0\}$. The inverse
continuous coefficient maps of complexes are
\[
 \begin{gathered}
 \mathscr C_N\simeq
 \mathcal C_N\oplus[\mathscr B^0\xrightarrow{I}\mathscr B^0],\\
 P\longmapsto(\av P,P-\Delta\av P),\qquad
 (v,b^0)\longmapsto\Delta v+b^0.
 \end{gathered}
 \tag{HCD6}
\]
In degree one, $J(P-\Delta\av P)=x-x=0$ and its mean is
zero. In degree zero every value already belongs to $B_N$.
Both composites restore the original vector. Inclusion
carries these components to inclusion and identity,
respectively, proving the isomorphism of complexes.

The degree-minus-one map
\[
 h:\mathscr P^{\rm eq}\longrightarrow\mathscr B,\qquad
 hP=P-\Delta\av P
 \tag{HCD7}
\]
satisfies $dh+hd=I-\Delta\av$. In degree one this is
$dhP=P-\Delta\av P$; in degree zero it is
$hd b=b-\Delta\av b$. On the complementary complex it
is the identity from degree one to degree zero. Boundedness
follows from (HCD5) and the isometry of repetition.
It is not asserted orthogonal or isometric.

The common quotient map is
\[
 \mathfrak j:\mathscr P^{\rm eq}\longrightarrow E,\qquad
 P\longmapsto JP(\theta)\text{ a.e.}
 \tag{HCD8}
\]
It is onto by repeating a right inverse of $J$, and its
kernel is exactly $\mathscr B$. Inclusion is injective.
Therefore $H^0(\mathscr C_N)=0$ and $H^1(\mathscr C_N)=E$,
with the isomorphism induced by the original full $J$.
Repetition and averaging induce the identity on this $E$.

Multiplication by $S$ maps $P_N$ to $P_{N+1}$ and $B_N$
to $B_{N+1}$. Applied pointwise, it sends common image $x$
to common image $Ax$ and commutes with repetition, mean
and $h$. These are maps between the stated cutoffs.
On degree-one cohomology the action is the full companion
$A$, including its nilpotent primary parts. No trace of
the infinite-dimensional contractible complement is assumed.

\section{The minimum norm of the common image}

Retain the original and phasewise minimum sections
\[
 \begin{aligned}
 G&=(JM^{-1}J^*)^{-1},& C&=M^{-1}J^*G,\\
 G_\theta&=(JM_\theta^{-1}J^*)^{-1},&
 C_\theta&=M_\theta^{-1}J^*G_\theta,\qquad
 \overline G=\int G_\theta\,d\mu.
 \end{aligned}
 \tag{HCD9}
\]
Surjectivity and positivity make every inverse defined.
Continuity of matrix inversion proves measurability.
For any fixed coefficient right inverse $C_{\rm ref}$ of $J$,
\[
 \frac{a}{\|J\|_{\rm coeff}^2}I_E
 \preceq G_\theta
 \preceq b\|C_{\rm ref}\|_{\rm coeff}^2I_E.
 \tag{HCD10}
\]
Indeed any lift $p$ of $x$ has
$\|x\|\le\|J\|\|p\|$, proving the lower bound, and the
particular lift $C_{\rm ref}x$ proves the upper bound.
The formula for $C_\theta$ then also gives boundedness.
All following integrals and sections are consequently defined.

Direct multiplication gives $JC_\theta=I$, and for $b_0\in B_N$,
\[
 b_0^*M_\theta C_\theta=b_0^*J^*G_\theta=0.
 \tag{HCD11}
\]
The section $\mathscr Sx(\theta)=C_\theta x$ has common image
$x$. Any other $P$ with that image has
$P-\mathscr Sx\in\mathscr B$. Phasewise orthogonality gives
\[
 \|P\|_{\mathscr P}^2
 =x^*\overline Gx+\|P-\mathscr Sx\|_{\mathscr P}^2.
 \tag{HCD12}
\]
Thus the quotient metric of the continuous complex is exactly
$\overline G>0$ on its full $E$. For the constant-source
complex the identical argument gives $G$. Repetition is
isometric on sources, and its induced identity
$(E,G)\to(E,\overline G)$ has the additional continuous
relation directions available in the target.

\section{The positive relation gap and its mean-zero decomposition}

For $r>0$, put
$X_\theta=(B^*M_\theta B)^{-1}B^*M_\theta C$.
Both sections have right-inverse image $I_E$, so their difference
lies in $B_N$ and equals $BX_\theta$; multiplication by
$B^*M_\theta$ proves the formula and uniqueness.
At $r=0$, both sections equal the unique inverse of $J$ and
$X_\theta$ is the zero map into the zero space.
Expanding $C=C_\theta+BX_\theta$ and using (HCD11) proves
\[
 \boxed{
 \begin{aligned}
 G-\overline G=\mathscr D
 &=\int(C-C_\theta)^*M_\theta(C-C_\theta)\,d\mu\\
 &=\int X_\theta^*(B^*M_\theta B)X_\theta\,d\mu\succeq0.
 \end{aligned}}
 \tag{HCD13}
\]
The integral of $C^*M_\theta C$ is $C^*MC=G$.
Every term is the norm of a field in the original relation
image. Multiplication by the literal $\chi$, followed by
the inherited finite tensor theta primitive, realizes
that field in the original cochain complex.

The relation to the particular mean-zero contraction is also exact.
Define
\[
 \overline C=\int C_\theta\,d\mu,\quad
 a_C=C-\overline C:E\to B_N,\quad
 h_\theta=C_\theta-\overline C:E\to B_N.
 \tag{HCD14}
\]
Then $\int h_\theta\,d\mu=0$, and $h\mathscr S$ is precisely
the field $h_\theta$. The original discrepancy is
$C-C_\theta=a_C-h_\theta$ and need not have mean zero.
Since $a_C$ has relation-valued columns, phase and
common-source orthogonality give
\[
 \begin{split}
 \int a_C^*M_\theta h_\theta\,d\mu
 &=-a_C^*M\overline C\\
 &=-a_C^*M(C-a_C)=a_C^*Ma_C.
 \end{split}
 \tag{HCD15}
\]
The adjoint cross term is the same Hermitian matrix.
Expanding all terms of (HCD13) therefore proves
\[
 \boxed{\mathscr D
 =\int h_\theta^*M_\theta h_\theta\,d\mu-a_C^*Ma_C.}
 \tag{HCD16}
\]
This retains the positive relation norm and the complete
weighted cross terms. A mean-zero coefficient field is
not presumed orthogonal to constant relations for the
varying metric. At $N=q-1$, $B_N=0$,
$C_\theta=C=J^{-1}$, $a_C=h_\theta=0$ and $\overline G=G$.

\section{Actual sampled Hilbert spaces, including zero weights}

Fix the original circumference $L>0$, tensor degree $k\ge1$,
and $c=k/2$. Retain
\[
 \begin{gathered}
 u_n(\theta)=(2\pi n+\theta)/L,\qquad
 s_n(\theta)=c+iu_n(\theta),\\
 w_n(\theta)=\frac{2\pi}{L}m_{h,k}(u_n(\theta))\ge0,
 \qquad n\in\mathbb Z.
 \end{gathered}
 \tag{HCD17}
\]
The Fourier coefficient is literally $2\pi/L$, with
$m_{h,k}=w_h^{*k}$, $w_h=|(g/h)(1/2+iu)|^2/(2\pi)$
and $g=2\xi$. Sampling gives
\[
 \mathsf T_\theta P=(P(s_n(\theta)))_n,\qquad
 \|\mathsf T_\theta P\|_\theta^2
 =\sum_n w_n(\theta)|P(s_n(\theta))|^2=P^*M_\theta P.
 \tag{HCD18}
\]
This is the actual phase source Gram H17.

Define $\mathcal H_\theta$ to be weighted square-summable
sequences modulo sequences supported where $w_n(\theta)=0$.
Thus zero weights, possible when $k=1$, have an explicit
null subspace. For $k\ge2$, positivity of the original
convolution makes every sample weight positive.
At $k=1$, retain the exact ideal equality
$(\chi_{h,1})=(h)$. If $a_h\ne0$ is the original leading
coefficient of $h$, then $\chi_{h,1}=h/a_h$ is its monic
generator, and $\chi_{h,1}=h$ in the original monic packet
convention. Indeed a polynomial annihilates multiplication by
$S$ on $\C[S]/(h)$ if and only if its action on $[1]$ is
zero, which is equivalent to belonging to $(h)$.
The analytic source remains the original $g/h$, with
the factor $a_h$ retained in that quotient and its norm.
Now take a complete selected root
$\rho=1/2+iu_n(\theta)$ whose common order in $g$ and $h$
is $m$. The two Taylor expansions have leading coefficients
$g^{(m)}(\rho)/m!$ and $h^{(m)}(\rho)/m!$, so the canceled
entire quotient has the exact nonzero value
\[
 (g/h)(\rho)=\frac{g^{(m)}(\rho)}{h^{(m)}(\rho)},\qquad
 w_n(\theta)=\frac1L
 \left|\frac{g^{(m)}(\rho)}{h^{(m)}(\rho)}\right|^2>0.
 \tag{HCD18a}
\]
This retains both cancellation factorials and the full
Fourier mass $(2\pi/L)(1/(2\pi))$. It is the selected-root
calculation PSD5 and PSD41. Other $k=1$ coordinates where
$g/h$ vanishes still give the explicit null subspace above.
For a concrete measurable realization use the isometry
$v\mapsto(\sqrt{w_n(\theta)}v_n)_n$ into ordinary
$\ell^2(\mathbb Z)$, with zero coordinates at zero weights.
Its inverse on its range is
$v_n=y_n/\sqrt{w_n(\theta)}$ at positive weights; other
coordinates give the null class. This specifies the norm
map without changing a mass. The range projections are
measurable diagonal matrices; their images of the countable
standard basis give a measurable field.

The original gamma-contour and alternating-zeta calculation
PSD8--PSD13 proves, for every $0<\gamma<\pi/2$, the finite
constants
\[
 A_\gamma=\|e^{\gamma|\cdot|}w_h\|_\infty,\qquad
 B_\gamma=\|e^{\gamma|\cdot|}w_h\|_1.
\]
That calculation retains the leading coefficient of $h$,
its full degree, and the continuous canceled values on the
compact interval containing its selected roots. With
$f_\gamma(v)=e^{\gamma|v|}w_h(v)$, the ordinary triangle
inequality on the convolution variables and integration
of the remaining factors give, for every real $u$,
\[
 e^{\gamma|u|}m_{h,k}(u)
 \le f_\gamma^{*k}(u)
 \le A_\gamma B_\gamma^{k-1}.
 \tag{HCD19a}
\]
For the second inequality bound one factor by $A_\gamma$
and integrate each of the other $k-1$ factors with its
literal Lebesgue measure. For $k=1$ the expression is
$A_\gamma$, with $B_\gamma^0=1$.
For any $0<2\epsilon<\gamma$,
\[
 \sup_{0\le\theta<2\pi}
 \sum_nw_n(\theta)e^{2\epsilon|u_n(\theta)|}<\infty.
 \tag{HCD19}
\]
Indeed $|u_n(\theta)|\ge(2\pi/L)(|n|-1)$, so with
$d=\gamma-2\epsilon>0$ the sum is at most
\[
 \frac{2\pi}{L}A_\gamma B_\gamma^{k-1}
 e^{2\pi d/L}
 \left(1+\frac{2e^{-2\pi d/L}}{1-e^{-2\pi d/L}}\right).
\]
For every fixed polynomial $Q$, the constant
$\sup_{u\in\R}|Q(c+iu)|^2e^{-d|u|/2}$ is finite.
Inserting it into the same geometric estimate, with
$d/2$ in place of $d$, proves uniform exponential
summability of $w_n(\theta)|Q(s_n(\theta))|^2$.
This includes $Q=\chi$ and every fixed polynomial
relation used below. All original constants in the
weights are unchanged. The full source proof of
PSD8--PSD14 is a dependency of this continuation;
the remaining discrete density argument is given
explicitly here with its shifted Fourier phase.

\section{Density and the complete fibrewise Hilbert quotient}

Here is a direct density proof with the zero weights retained.
If $f\in\mathcal H_\theta$ is orthogonal to all polynomials,
put $a_n=w_n(\theta)\overline{f_n}$. Cauchy--Schwarz and
(HCD19) give $\sum|a_n|e^{\epsilon|u_n|}<\infty$.
Hence $F(z)=\sum_n a_ne^{zu_n}$ is holomorphic on
$|\Re z|<\epsilon$, with termwise derivatives on smaller
closed strips. Every derivative at zero is zero, since
polynomials in $S$ and in $(S-c)/i$ span the same space.
The Taylor series and the identity theorem give $F=0$
on the connected strip. On the imaginary axis,
\[
 F(it)=e^{it\theta/L}\sum_n a_ne^{2\pi int/L}.
\]
The series is absolutely uniformly convergent.
Multiplication by $e^{-it\theta/L}e^{-2\pi imt/L}$ and
integration over $0\le t\le L$ give $a_m=0$.
At positive weights this implies $f_m=0$; at zero weights
the class is already null. Thus polynomials are dense.
The identical argument proves polynomial density for
the weights $w_n|\chi(s_n)|^2$.

Multiplication by the original polynomial is the isometry
\[
 U_\chi:\ell^2(w_n|\chi(s_n)|^2)\longrightarrow\mathcal H_\theta,
 \qquad f\longmapsto(\chi(s_n)f_n)_n.
 \tag{HCD20}
\]
Its image is precisely the sequences vanishing at the
positive-weight sampled roots of $\chi$. To prove the
converse inclusion, take such a sequence $v$ and set
$f_n=v_n/\chi(s_n)$ when $\chi(s_n)\ne0$, and zero
elsewhere. Its squared domain norm equals $\|v\|_\theta^2$.
Thus the image is closed. Polynomial density in the domain
now proves
\[
 \mathcal R_\theta:=
 \overline{\chi\C[S]}^{\,\mathcal H_\theta}
 =\{v:v_n=0\text{ if }w_n(\theta)>0,\ \chi(s_n(\theta))=0\}.
 \tag{HCD21}
\]
The receiving Hilbert quotient is exactly
\[
 \begin{gathered}
 \mathcal Q_\theta=\mathcal H_\theta/\mathcal R_\theta
 \simeq
 \bigoplus_{\substack{n:\chi(s_n(\theta))=0\\w_n(\theta)>0}}\C,\\
 \|v\|_{\mathcal Q_\theta}^2
 =\sum_{\rm displayed\ n}w_n(\theta)|v_n|^2.
 \end{gathered}
 \tag{HCD22}
\]
The quotient map is orthogonal coordinate restriction.
No weight equal to zero is inverted.

\section{The complete primary kernel before phase integration}

Let $\Lambda_\theta^+$ be the distinct sampled roots of
$\chi$ with positive weight, and set
$p_\theta(S)=\prod_{\lambda\in\Lambda_\theta^+}(S-\lambda)$,
using $p_\theta=1$ for an empty product.
The finite map is
\[
 \rho_\theta:E\longrightarrow\mathcal Q_\theta,\qquad
 [P]_\chi\longmapsto(P(\lambda))_{\lambda\in\Lambda_\theta^+}.
 \tag{HCD23}
\]
It is well defined, and is onto by Lagrange interpolation
at these distinct points, using degree less than
$|\Lambda_\theta^+|\le q$. Vanishing at all those points
is equivalent to divisibility by $p_\theta$. Since
$p_\theta\mid\chi$,
\[
 \ker\rho_\theta=(p_\theta)/(\chi)\subset E.
 \tag{HCD24}
\]
This specifies the zero-weight-corrected kernel in H64:
a root at a zero-mass coordinate contributes no receiving
Hilbert line.

Write the full original factorization as
$\chi=\prod_{\lambda\in\Lambda}(S-\lambda)^{m_\lambda}$.
The complete Taylor map is
\[
 \begin{split}
 E&\longrightarrow\prod_{\lambda\in\Lambda}
 \C[\varepsilon_\lambda]/(\varepsilon_\lambda^{m_\lambda}),\\
 [P]&\longmapsto\left(
 \sum_{j=0}^{m_\lambda-1}\frac{P^{(j)}(\lambda)}{j!}
 \varepsilon_\lambda^j\right)_\lambda.
 \end{split}
 \tag{HCD25}
\]
For an explicit inverse, let
$q_\lambda=\chi/(S-\lambda)^{m_\lambda}$ and let
$u_\lambda$ be the Taylor polynomial of $1/q_\lambda$
at $\lambda$ through degree $m_\lambda-1$.
Then $e_\lambda=[u_\lambda q_\lambda]_\chi$ is one
in the $\lambda$ factor and zero in every other factor,
by divisibility by each original prime power. It follows
that $\sum e_\lambda=1$ and the factors are orthogonal
idempotents. The inverse is
$\sum_\lambda e_\lambda\sum_j a_{\lambda,j}(S-\lambda)^j$.

In these full primary coordinates, $\rho_\theta$ takes
the constant coefficient of each active sampled factor.
Its kernel in that factor is $(\varepsilon_\lambda)$,
retaining every order $1,\ldots,m_\lambda-1$.
For an unsampled factor or a zero-weight sampled factor,
the kernel is the entire primary algebra, including
its constant term. The original action is multiplication
by $\lambda+\varepsilon_\lambda$; the receiving action
is the simple scalar $\lambda$. Evaluation intertwines
these actions with the displayed nilpotent kernel.

\section{The measurable completed field and its vanishing map}

Every sampled root has real part $c$. For a fixed root $\lambda$
with $\Re\lambda=c$, its only possible phase is
\[
 \theta_\lambda=L\Im\lambda\pmod{2\pi},\quad
 0\le\theta_\lambda<2\pi,\qquad
 n_\lambda=(L\Im\lambda-\theta_\lambda)/(2\pi)\in\mathbb Z.
\]
Thus the set $\Theta_\chi$ of possible phases is finite.
For $\theta\notin\Theta_\chi$, $\mathcal Q_\theta=0$.
At an exceptional phase the positive-weight condition
in (HCD22) is still retained.

In the measurable realization after (HCD18), the projection
onto $\mathcal Q_\theta$ has diagonal entries
\[
 {\bf1}_{\{w_n(\theta)>0\}}
 {\bf1}_{\{\chi(s_n(\theta))=0\}},\qquad n\in\mathbb Z.
\]
These are measurable and give an explicit field of projections.
The direct integral satisfies
\[
 \mathscr Q=\int_\Theta^\oplus\mathcal Q_\theta\,d\mu=0,
 \tag{HCD26}
\]
because every measurable section has squared norm supported
on the finite measure-zero set $\Theta_\chi$.
This is an almost-everywhere assertion; each exceptional
nonzero receiving line is still exactly the weighted line
in (HCD22).

For $x\in E$, the section $\theta\mapsto\rho_\theta x$
is measurable: it is given by polynomial evaluations and
the displayed coordinate masks. It defines
\[
 \rho:E\longrightarrow\mathscr Q,\qquad
 x\longmapsto[\theta\mapsto\rho_\theta x]_{\rm a.e.},
 \qquad \ker\rho=E.
 \tag{HCD27}
\]
Every pointwise exceptional value has zero direct-integral
class. This does not assert that those pointwise maps vanish,
or that the original nonreduced coefficient algebra was zero.

\section{The source-to-completion cochain map}

The direct-integral source is the weighted space with squared
norm $\int_\Theta\sum_n w_n(\theta)|v_n(\theta)|^2\,d\mu$.
Let $\mathscr H$ denote this space, and let $\mathscr R$
be its closed subspace of sections belonging to
$\mathcal R_\theta$ almost everywhere. The coordinate
projection above has kernel $\mathscr R$.

The measurable evaluation map is
\[
 \mathsf T:\mathscr P\longrightarrow\mathscr H,\qquad
 P\longmapsto(P(\theta,s_n(\theta)))_{\theta,n}.
\]
The integrated identity (HCD18) proves it is an isometry.
Evaluation of a relation-valued section is in $\mathscr R$,
because the original section is a multiple of the literal
$\chi$. Hence there is a cochain map
\[
 [\mathscr B\hookrightarrow\mathscr P^{\rm eq}]
 \longrightarrow[\mathscr R\hookrightarrow\mathscr H],
 \tag{HCD28}
\]
using evaluation in both degrees. On a common image $x$,
the receiving quotient value is $\rho_\theta x$, since
representatives differ by an original relation.
The induced degree-one map is exactly (HCD27).

The global relation subspace is also the closure of measurable
finite polynomial relation sections. To prove this without
a measurable choice of approximating polynomials, suppose
$v\in\mathscr H$ is orthogonal to every section
$\mathbf1_U(\theta)\chi(S)S^j$, for measurable $U$ and $j\ge0$.
Its fibre inner product with $\chi S^j$ is integrable by
Cauchy--Schwarz, and its integral over every $U$ is zero,
so that inner product is zero almost everywhere.
The original exponential sample bounds make every fixed
polynomial section square-integrable.
Intersecting the countably many resulting full-measure sets,
(HCD21) gives $v(\theta)\in\mathcal Q_\theta$ there.
Conversely every such section is orthogonal to the relations.
The orthogonal-complement characterization proves that their
closed span is exactly $\mathscr R$.

Since $\mathscr Q=0$, $\mathscr R=\mathscr H$ as Hilbert
spaces, and the completed complex is the identity complex.
Its degree-one cohomology is zero. The finite continuous
common-image complex has positive Gram $\overline G$
on its full $E$ by (HCD8)--(HCD12).
The intervening map and its whole kernel are precisely
(HCD23)--(HCD28).

\section{The receiving supported zero and all marked fibres}

Integrate in the stated coefficient fibre, then apply the
original marked support construction. At a fixed original
label $\ell$, a linear map $f$ has the lift
\[
 (\ell,x)\longmapsto(\ell,f(x)),\qquad
 \tau\longmapsto\tau.
 \tag{HCD29}
\]
For (HCD27), every represented input maps to the receiving
supported zero $(\ell,0)$, and external absence maps to
external absence $\tau$. Their inverse images are respectively
the full represented fibre and the singleton external absence.
At a fixed phase, the represented kernel is the entire
ideal (HCD24), with its full primary decomposition (HCD25).
No integral on arbitrary mixed-support totals is presumed.

If the empty-packet case $q=0$ is separately admitted,
$E=0$, $B_N=P_N$, and the finite and continuous coefficient
complexes are identity complexes. The quotient Gram
determinant is the empty determinant one. The analytic
source and its literal norm remain present before the
quotient. Its marked receiving object still has $\tau$
and the represented zero, with the same maps (HCD29).

\section{Exact outcome and scope}

The continuous common-image complex has the explicit bounded
cochain contraction (HCD7) on its complement and retains
the whole nonreduced quotient in degree one.
Its minimum metric is $\overline G$, and its difference
from the constant-source metric is the positive original
relation norm (HCD13). The mean-zero decomposition retains
the weighted cross term (HCD15), giving (HCD16).

The subsequent completion has receiving weighted value
spaces (HCD22), primary kernels (HCD24)--(HCD25), measurable
field (HCD26), and actual cochain map (HCD28).
Its receiving coefficient direct integral is zero because
only finitely many phases can sample a fixed polynomial root.
Every exceptional phase, nilpotent component and zero-weight
coordinate remains in the stated maps and kernels.
No arithmetic volume asymptotic or RH conclusion is inferred.

